\documentclass[12pt,letterpaper]{article}

% Paquetes esenciales
\usepackage[utf8]{inputenc}
\usepackage[spanish,es-noquoting]{babel} 
\usepackage{geometry}
\geometry{top=2.5cm, bottom=2.5cm, left=3cm, right=3cm}
\usepackage{tikz}
\usetikzlibrary{shapes.geometric, arrows.meta, positioning, fit, backgrounds, calc}
\usepackage{enumitem}
\usepackage{titlesec}
\usepackage{hyperref}
\usepackage{booktabs}
\usepackage{array}
\usepackage{graphicx}

% Formato de títulos
\titleformat{\section}{\large\bfseries}{\thesection}{1em}{}
\titleformat{\subsection}{\normalsize\bfseries}{\thesubsection}{1em}{}

\begin{document}

% ==========================================
% PORTADA
% ==========================================
\begin{center}
    \vspace*{1cm}
    \Large \textbf{Propuesta de Proyecto: Traductor de Lengua de Señas Colombiana (LSC) mediante Arquitectura YOLO} \\
    \vspace{0.5cm}
    \normalsize \textbf{Tomas Patiño} \\
    \vspace{0.2cm}
    \textit{Introducción a la Visión por Computador} \\
    \vspace{1cm}
\end{center}

% ==========================================
% 1. OBJETIVOS
% ==========================================
\section{Objetivos}

\subsection{Objetivo General}
Desarrollar un sistema de visión artificial basado en aprendizaje profundo para la detección y clasificación en tiempo real del alfabeto estático de la Lengua de Señas Colombiana (LSC), utilizando una arquitectura de detección de objetos de un solo paso (\textit{End-to-End}).

\subsection{Objetivos Específicos}
\begin{itemize}[leftmargin=*]
    \item Construir y consolidar un conjunto de datos (\textit{dataset}) propio y representativo mediante la captura fotográfica de 10 voluntarios, garantizando variabilidad en la fisonomía, tonos de piel y fondos.
    \item Implementar un flujo de pre-procesamiento que incluya la estandarización espacial de las imágenes y la anotación manual de características mediante cajas delimitadoras (\textit{bounding boxes}).
    \item Entrenar un modelo basado en la arquitectura YOLO, delegando la extracción de características visuales a su red convolucional base (\textit{Backbone}).
    \item Evaluar el rendimiento del modelo simultáneo de localización y clasificación utilizando la métrica de Precisión Media Promedio (mAP) en entornos no controlados.
\end{itemize}

% ==========================================
% 2. ESTADO DEL ARTE
% ==========================================
\section{Investigación del Estado del Arte}
La literatura reciente (2023-2025) evidencia una transición desde modelos de clasificación por etapas hacia arquitecturas unificadas para el reconocimiento de lenguaje de señas. 

Tyagi et al. (2023) demostraron que la arquitectura YOLO supera a las Redes Neuronales Convolucionales (CNN) tradicionales en la detección del alfabeto de señas (ASL), logrando un mAP del 96\% en tiempo real, incluso bajo limitaciones de hardware. Por otro lado, investigaciones recientes como la de Madyanto y Wijaya (2025) validan la eficacia de modelos YOLO (específicamente YOLOv8) para la detección de manos en entornos no controlados con fondos complejos, alcanzando precisiones superiores al 90\% frente a variaciones severas de iluminación. Finalmente, diversos estudios comparativos recientes (2025) subrayan que el rendimiento superior al 97\% en estos sistemas depende inherentemente de una estrategia rigurosa de \textit{Data Augmentation} para prevenir el sobreajuste.

% ==========================================
% 3. APLICACIÓN DE LA DOCUMENTACIÓN
% ==========================================
\section{Aplicación de la Documentación Investigada}
La evidencia teórica recolectada fundamenta las decisiones estructurales del presente proyecto:
\begin{itemize}[leftmargin=*]
    \item \textbf{Selección de la Arquitectura:} En respuesta a los hallazgos de Tyagi et al. (2023), se descarta el uso de clasificadores de imágenes estáticos o arquitecturas de múltiples etapas (recorte manual seguido de clasificación). Se adopta YOLO para garantizar un procesamiento \textit{End-to-End}.
    \item \textbf{Diseño del Dataset:} Basado en el estudio de Madyanto y Wijaya (2025) sobre fondos complejos, se justifica la creación de un dataset propio con 10 personas distintas. Esto entrenará al modelo para ignorar el ruido visual del fondo y centrarse estrictamente en la morfología de la seña.
    \item \textbf{Estrategia de Pre-procesamiento:} Siguiendo las directrices metodológicas de 2025 sobre aumento de datos, se aplicarán rotaciones leves y ajustes fotométricos, omitiendo estrictamente el efecto espejo (\textit{horizontal flip}) para preservar el significado lateral de la LSC. Asimismo, se delega la extracción de características al \textit{Backbone} de YOLO, eliminando el uso de filtros manuales de visión clásica.
\end{itemize}

% ==========================================
% 4. METODOLOGÍA
% ==========================================
\section{Metodología}
El desarrollo del proyecto se ejecutará mediante un enfoque secuencial estructurado en tres capas lógicas principales:

\textbf{1. Capa de Pre-procesamiento y Aumento de Datos (Input Layer):}
Inicia con la recolección de las fotografías del alfabeto LSC. Las imágenes estáticas serán consolidadas y sometidas a un proceso de normalización, redimensionándolas a matrices cuadradas de $640 \times 640$ píxeles. Previo al entrenamiento, se aplicarán técnicas de \textbf{Data Augmentation} (aumento de datos) controladas para robustecer el modelo, incluyendo ajustes fotométricos (brillo y contraste) y variaciones geométricas leves (rotaciones de $\pm 15^\circ$), excluyendo estrictamente la inversión lateral (\textit{horizontal flip}) para no alterar el significado direccional de las señas. Posteriormente, se realizará la anotación espacial manual trazando las \textit{bounding boxes} y asignando la etiqueta correspondiente a cada letra. 

\textbf{2. Capa de Modelado (Shared Layers / Backbone):}
El conjunto de datos anotado se dividirá en proporciones para entrenamiento, validación y prueba. Las matrices de entrada se procesarán mediante la red convolucional base (\textit{YOLO Backbone}), la cual extraerá los mapas de características jerárquicas sin requerir filtros de extracción de bordes externos, optimizando los núcleos de convolución durante el entrenamiento.

\textbf{3. Capa de Predicción (Task Layers / YOLO Head):}
Las características extraídas se canalizarán hacia las cabezas de predicción, las cuales ejecutarán tres tareas simultáneas: la regresión de las coordenadas espaciales de la mano, el cálculo de la probabilidad de presencia del objeto y la clasificación probabilística de la letra. Finalmente, se aplicará el algoritmo de Supresión de No Máximos (NMS) para unificar predicciones redundantes y se evaluará el sistema usando la métrica mAP.

\vspace{0.5cm}

% ==========================================
% 5. DESCRIPCIÓN DEL PREPROCESAMIENTO
% ==========================================
\section{Descripción del preprocesamiento de los datos}

Para la consolidación del conjunto de datos (\textit{dataset}) propio, se diseñó e implementó una canalización (\textit{pipeline}) de captura automatizada mediante el uso de librerías de visión por computador (OpenCV). Esta decisión metodológica se tomó para garantizar un control estricto sobre la variabilidad y el balance de las clases, evitando depender de bases de datos genéricas que no se ajustan a las particularidades morfológicas y de iluminación del entorno real de aplicación.

El proceso de recolección se estructuró bajo los siguientes parámetros técnicos:
\begin{itemize}
    \item \textbf{Gestión de Voluntarios y Clases:} Se estableció un protocolo de captura para las 21 letras estáticas del alfabeto de la Lengua de Señas Colombiana (LSC). Con 10 voluntarios. El sistema gestiona las sesiones mediante identificadores únicos para cada voluntario (por ejemplo, V01, V02), asegurando la recolección de muestras equitativas por sujeto y garantizando la diversidad en tonos de piel y proporciones anatómicas.
    \item \textbf{Estructuración Topológica:} Las imágenes capturadas son almacenadas automáticamente en una jerarquía de directorios estandarizada (\texttt{dataset/\{Letra\}/\{Voluntario\}/}). Esto facilita la posterior partición iterativa de los datos para las fases de entrenamiento, validación y prueba.
    \item \textbf{Trazabilidad y Metadatos:} De forma concurrente a la captura visual, el sistema genera registros estructurados en formato JSON. Estos bitácoras documentan la marca de tiempo (\textit{timestamp}), el identificador del sujeto y la clase correspondiente, asegurando la reproducibilidad del experimento.
    \item \textbf{Preservación de la Direccionalidad:} Durante la captura, se implementó un efecto de espejo (\textit{horizontal flip}) exclusivamente en la previsualización de la pantalla para mejorar la ergonomía del voluntario. Sin embargo, la matriz de píxeles se almacena en su orientación original, cumpliendo con la restricción de no alterar la direccionalidad lateral fundamental de las señas en la LSC.
\end{itemize}

% ==========================================
% 6. PRIMEROS ACERCAMIENTOS
% ==========================================
\section{Primeros acercamientos al desarrollo del modelo}

En las fases iniciales del desarrollo, el enfoque metodológico consistió en abordar el problema como una tarea clásica de Detección de Objetos. Para ello, se estructuró una herramienta que delimitaba la región de interés (ROI) mediante cajas delimitadoras ortogonales (\textit{Bounding Boxes}). Sin embargo, el análisis preliminar de estas muestras reveló que este método encapsulaba un alto porcentaje de ``ruido visual'' (píxeles pertenecientes al fondo), lo cual representa un riesgo significativo para la precisión de la clasificación, dado que en el lenguaje de señas la forma exacta de los dedos es el único discriminante válido.

Ante este desafío, se tomó la decisión estratégica de evolucionar la arquitectura del proyecto hacia una tarea de \textbf{Segmentación de Instancias}. Este pivote tecnológico permite que el modelo aísle computacionalmente la mano, ignorando por completo el entorno no controlado.

Para materializar este acercamiento, se implementaron las siguientes acciones:
\begin{itemize}
    \item \textbf{Anotación Semántica por Polígonos:} Se integró la herramienta de anotación \textit{Labelme} al flujo de trabajo. Mediante esta plataforma, se están trazando máscaras poligonales vértice a vértice que contornean con exactitud la silueta de la mano en cada fotografía capturada.
    \item \textbf{Preparación para YOLOv8-seg:} El trazado manual de polígonos exporta las coordenadas en formato estructurado (JSON). Como paso de integración, se configuró un entorno virtual aislado con librerías de conversión (\texttt{labelme2yolo}), las cuales transforman estas máscaras poligonales al formato de tensores normalizados exigido por la variante de segmentación del modelo YOLO (YOLOv8-seg).
\end{itemize}

Esta evolución en el desarrollo preliminar garantiza que la red neuronal convolucional enfoque la actualización de sus pesos matemáticos exclusivamente en las características morfológicas y geométricas de las señas, maximizando la robustez del sistema frente a fondos complejos.

% ==========================================
% 7. PROCEDIMIENTOS REALIZADOS EN EL PROYECTO
% ==========================================
\section{Procedimientos Realizados en el Proyecto}

La implementación del sistema siguió un \textit{pipeline} secuencial de siete etapas, cada una materializada en un módulo de software independiente.

\subsection{Captura del Dataset}
Se desarrolló una herramienta de captura interactiva (\texttt{capturar\_dataset.py}) utilizando OpenCV. La interfaz muestra en tiempo real la letra objetivo, el progreso por clase y una barra de avance. El sistema captura un frame limpio (sin \textit{overlays} de interfaz) al momento de la foto, lo almacena bajo la jerarquía \texttt{dataset/\{Letra\}/\{Voluntario\}/} y registra los metadatos en un archivo \texttt{log\_captura.json}.

\subsection{Etiquetado con LabelMe}
Cada imagen fue anotada manualmente en LabelMe trazando un \textbf{polígono} sobre la silueta completa de la mano. Se optó por polígonos en lugar de \textit{bounding boxes} para capturar con exactitud la forma de los dedos, único discriminante visual entre muchas señas. El etiquetado generó un archivo \texttt{.json} por imagen con las coordenadas absolutas de los vértices y la etiqueta de clase.

\subsection{Conversión de Formato (\texttt{labelme2yolo\_converter.py})}
Se implementó un conversor que transforma los archivos \texttt{.json} de LabelMe al formato \texttt{.txt} exigido por YOLOv8 Segmentation:
\[
\texttt{<class\_id>} \quad x_1 \; y_1 \quad x_2 \; y_2 \quad \cdots \quad x_n \; y_n
\]
donde todas las coordenadas están normalizadas en $[0,\,1]$. El script procesa el directorio \texttt{dataset/} de forma recursiva y reporta el conteo de conversiones exitosas y omitidas.

\subsection{Construcción de la Estructura YOLO (\texttt{build\_yolo\_structure.py})}
El dataset se reorganizó en la estructura canónica de YOLOv8 con una partición estratificada por clase (semilla fija $= 42$): \textbf{80\% entrenamiento}, \textbf{10\% validación} y \textbf{10\% prueba}. Se generó automáticamente el archivo \texttt{dataset.yaml} con las rutas absolutas y la lista de las 21 clases.

\subsection{Aumento de Datos (\texttt{data\_augmentation.py})}
Para compensar el tamaño reducido del dataset original, se multiplicó $\times 5$ el split de entrenamiento mediante la librería \textbf{Albumentations}, usando \texttt{KeypointParams} para transformar los polígonos de forma coherente con la imagen. Las transformaciones aplicadas se resumen en la Tabla~\ref{tab:augmentation}.

\begin{table}[h!]
\centering
\caption{Pipeline de \textit{Data Augmentation} aplicado al split de entrenamiento.}
\label{tab:augmentation}
\begin{tabular}{@{}lll@{}}
\toprule
\textbf{Transformación} & \textbf{Parámetros} & \textbf{Prob.} \\ \midrule
Flip horizontal         & ---                              & 50\%  \\
ShiftScaleRotate        & $\pm$5\% despl., $\pm$15\% escala, $\pm$20° rot. & 80\%  \\
RandomBrightnessContrast & $\pm$30\% brillo y contraste    & 70\%  \\
HueSaturationValue      & $\pm$15° hue, $\pm$30 sat, $\pm$20 val & 50\%  \\
GaussianBlur / MotionBlur & kernel 3--5 px               & 30\%  \\
GaussNoise              & varianza 5--30                   & 30\%  \\
CLAHE                   & clip 2.0, grilla 8$\times$8      & 20\%  \\ \bottomrule
\end{tabular}
\end{table}

\subsection{Entrenamiento con Transfer Learning (\texttt{train\_yolo.py})}
Se entrenó el modelo \textbf{YOLOv8n-seg} (variante \textit{nano}, $\approx$3.4M parámetros) partiendo de pesos COCO preentrenados. La GPU utilizada fue una NVIDIA RTX 2060 (6 GB VRAM). Los hiperparámetros principales se detallan en la Tabla~\ref{tab:hiper}.

\begin{table}[h!]
\centering
\caption{Hiperparámetros de entrenamiento del modelo final (v15).}
\label{tab:hiper}
\begin{tabular}{@{}lll@{}}
\toprule
\textbf{Parámetro} & \textbf{Valor} & \textbf{Justificación} \\ \midrule
Épocas máximas     & 150            & Con \textit{early stopping} (paciencia = 30) \\
Batch size         & 4              & Limitado por 6 GB VRAM \\
Resolución entrada & 640$\times$640 px & Estándar YOLOv8 \\
Optimizador        & AdamW          & Más estable para datasets pequeños \\
Learning rate (lr0)& 0.001          & Fine-tuning conservador \\
AMP (FP16)         & Activado       & Reduce VRAM $\approx$50\% \\
Pretrained         & Activado       & Transfer Learning desde COCO \\ \bottomrule
\end{tabular}
\end{table}

\subsection{Inferencia en Tiempo Real (\texttt{test\_webcam.py})}
El modelo entrenado (\texttt{best.pt}) se desplegó sobre el stream de la cámara web. Cada frame se espeja horizontalmente antes de la inferencia (consistente con la etapa de captura) y las predicciones se visualizan con las máscaras de segmentación superpuestas mediante \texttt{results[0].plot()} de Ultralytics. Se utilizó un umbral de confianza de 0.15 para maximizar la sensibilidad en la demostración.

% ==========================================
% 8. RESULTADOS
% ==========================================
\section{Resultados}

\subsection{Progresión del Entrenamiento}
El modelo fue entrenado durante \textbf{132 épocas efectivas}, alcanzando la convergencia antes del límite máximo de 150. La Tabla~\ref{tab:progresion} muestra la evolución de las métricas clave en puntos representativos del entrenamiento.

\begin{table}[h!]
\centering
\caption{Evolución de métricas en el conjunto de validación durante el entrenamiento.}
\label{tab:progresion}
\begin{tabular}{@{}cccc@{}}
\toprule
\textbf{Época} & \textbf{mAP50 (B)} & \textbf{mAP50 (M)} & \textbf{mAP50-95 (M)} \\ \midrule
10  & 75.2\% & 75.2\% & 63.7\% \\
30  & 93.0\% & 93.0\% & 83.4\% \\
50  & 96.6\% & 96.6\% & 86.4\% \\
70  & 99.5\% & 99.5\% & 89.2\% \\
100 & 99.5\% & 99.5\% & 90.3\% \\
132 & \textbf{99.5\%} & \textbf{99.5\%} & \textbf{91.5\%} \\ \bottomrule
\end{tabular}
\end{table}

\subsection{Métricas Finales del Modelo}
La Tabla~\ref{tab:metricas} presenta las métricas del mejor modelo obtenido (experimento \textit{senas\_colombianas\_v15}), evaluado sobre el conjunto de validación al finalizar el entrenamiento.

\begin{table}[h!]
\centering
\caption{Métricas finales del modelo YOLOv8n-seg entrenado (v15).}
\label{tab:metricas}
\begin{tabular}{@{}lc@{}}
\toprule
\textbf{Métrica} & \textbf{Valor} \\ \midrule
mAP50 — \textit{Bounding Box}             & \textbf{99.5\%} \\
mAP50 — Máscara de Segmentación           & \textbf{99.5\%} \\
mAP50-95 — Máscara de Segmentación        & \textbf{91.5\%} \\
Precisión (Bounding Box)                  & 90.5\%  \\
Recall (Bounding Box)                     & 99.5\%  \\
Pérdida de clasificación (validación)     & 0.433   \\
Pérdida de \textit{box} (validación)      & 0.424   \\
Pérdida de segmentación (validación)      & 0.571   \\ \bottomrule
\end{tabular}
\end{table}

El mAP50 de 99.5\% indica que el modelo prácticamente no comete errores de clasificación al umbral $\text{IoU} = 0.50$ en el conjunto de validación. El mAP50-95 de 91.5\% es la métrica más exigente (promedio sobre umbrales IoU de 0.50 a 0.95) y refleja la calidad geométrica de las máscaras generadas. Adicionalmente, el modelo final pesa aproximadamente \textbf{5.5 MB}, lo que lo hace apto para despliegue en dispositivos de recursos limitados.

\subsection{Rendimiento en Tiempo Real}
Sobre la GPU NVIDIA RTX 2060, el sistema realiza inferencias a más de \textbf{30 fps} sobre el stream de la cámara web, cumpliendo el requisito de operación en tiempo real. Las máscaras de segmentación permiten distinguir señas visualmente similares (p.~ej., U/V, M/N, K/P) con alta fidelidad.

% ==========================================
% 9. CONCLUSIONES
% ==========================================
\section{Conclusiones}

\begin{enumerate}[leftmargin=*, label=\arabic*.]
    \item \textbf{Viabilidad del enfoque:} Es posible construir un sistema de reconocimiento de señas LSC de alta precisión (mAP50 $>$ 99\%) usando únicamente una cámara web estándar y herramientas de código abierto, sin requerir guantes sensoriales ni cámaras de profundidad.

    \item \textbf{Efectividad del \textit{transfer learning}:} Los pesos COCO preentrenados proporcionaron una inicialización que permitió convergencia rápida incluso con un dataset pequeño ($\approx$300 imágenes en entrenamiento tras augmentation), validando que las características aprendidas para detección de objetos generales son transferibles al reconocimiento de gestos de mano.

    \item \textbf{Relevancia del aumento de datos:} La multiplicación $\times$5 del dataset de entrenamiento mediante Albumentations fue determinante para la generalización. Sin augmentation, el riesgo de sobreajuste habría sido alto dado el tamaño original del dataset.

    \item \textbf{Superioridad de la segmentación de instancias:} El pivote de \textit{bounding boxes} a polígonos de segmentación aportó valor real al aislar la silueta exacta de la mano, lo cual es crítico para señas que solo se diferencian en la posición de los dedos (p.~ej., M vs.~N, U vs.~V).

    \item \textbf{Limitaciones identificadas:} El sistema muestra sensibilidad a condiciones de iluminación extremas y fondos muy cambiantes, ya que el dataset de entrenamiento fue capturado en condiciones relativamente controladas. Ampliar la diversidad de voluntarios y entornos mejoraría la robustez del modelo.

    \item \textbf{Trabajo futuro:} Los pasos naturales de continuación incluyen: (a) extensión a señas dinámicas mediante análisis de secuencias de frames, (b) integración de un módulo de síntesis de texto para formar palabras completas a partir de señas reconocidas en secuencia, y (c) despliegue del sistema en una aplicación web o móvil de acceso público para impacto social directo.
\end{enumerate}

% ==========================================
% REFERENCIAS
% ==========================================
\section*{Referencias Bibliográficas}
\begin{itemize}[leftmargin=*]
    \item Tyagi, C., et al. (2023). \textit{American Sign Language Detection using YOLOv5 and YOLOv8}. Preprint / arXiv.
    \item Madyanto, R. K., \& Wijaya, Y. A. (2025). YOLOv8 Algorithm to Improve the Sign Language Letter Detection System Model. \textit{Journal of Artificial Intelligence and Engineering Applications (JAIEA)}, 4(2), 1379--1385.
    \item Buslaev, A., et al. (2020). Albumentations: Fast and Flexible Image Augmentations. \textit{Information}, 11(2), 125.
    \item Jocher, G. et al. (2023). \textit{Ultralytics YOLOv8}. \url{https://github.com/ultralytics/ultralytics}
\end{itemize}

\begin{figure}[htbp] 
    \centering
    \resizebox{\textwidth}{!}{
    \begin{tikzpicture}[
        node distance=1.5cm and 2cm,
        caja/.style={rectangle, rounded corners, draw=black, thick, text width=3.5cm, align=center, minimum height=1.2cm, fill=white},
        matriz/.style={rectangle, draw=black, dashed, text width=3cm, align=center, minimum height=1cm, fill=gray!10},
        flecha/.style={thick, ->, >=stealth},
        linea/.style={thick, -}
        ]

        % ==========================================
        % 1. INPUT LAYER (Pre-procesamiento)
        % ==========================================
        \node (db) [cylinder, draw=black, thick, aspect=0.25, minimum height=2cm, minimum width=1.5cm, shape border rotate=90, fill=black!70, text=white, align=center] {Dataset \\ LSC};
        
        \node (fotos) [below=0.2cm of db, text width=3cm, align=center, font=\scriptsize] {Fotos de 10 personas realizando el alfabeto};

        \node (ext) [caja, right=1.5cm of db] {Extracción y Normalización \\ (640x640px)};
        \node (sam) [caja, right=of ext] {Anotación Espacial \\ (\textit{Bounding Boxes})};
        \node (metric) [caja, right=of sam] {Etiquetado de Clases \\ (A - Z)};

        \draw [flecha] (db) -- (ext);
        \draw [flecha] (ext) -- (sam);
        \draw [flecha] (sam) -- (metric);

        \begin{scope}[on background layer]
            \node[draw=blue, thick, dotted, fill=blue!5, inner sep=15pt, fit=(ext) (sam) (metric)] (input_layer) {};
            \node[anchor=south east, text=blue, font=\bfseries\footnotesize] at (input_layer.south east) {Input Layer};
        \end{scope}

        % ==========================================
        % 2. SHARED LAYERS (Backbone YOLO)
        % ==========================================
        \node (split) [caja, below=2cm of sam] {División de Datos \\ (Train, Val, Test)};
        \node (backbone) [caja, below=1cm of split, fill=green!10] {YOLO Backbone \\ (Extracción de Features)};
        
        \node (tensor_in) [matriz, left=1cm of backbone] {Tensor Entrada \\ $\{ N \times 3 \times H \times W \}$};
        \node (tensor_out) [matriz, right=1cm of backbone] {Feature Maps \\ $\{ S \times S \times C \}$};
        \node (transfer) [caja, above=0.8cm of tensor_out, text width=2.5cm, minimum height=0.8cm] {Transfer Learning};

        \draw [flecha] (metric.south) |- (split.east);
        \draw [flecha] (split) -- (backbone);
        \draw [flecha] (tensor_in) -- (backbone);
        \draw [flecha] (backbone) -- (tensor_out);
        \draw [flecha] (transfer) -- (tensor_out);

        \begin{scope}[on background layer]
            \node[draw=orange, thick, dotted, fill=orange!5, inner sep=15pt, fit=(split) (backbone) (tensor_in) (tensor_out) (transfer)] (shared_layer) {};
            \node[anchor=south east, text=orange, font=\bfseries\footnotesize] at (shared_layer.south east) {Shared Layers (Backbone \& Neck)};
        \end{scope}

        % ==========================================
        % 3. TASK LAYERS (Predicción / YOLO Head)
        % ==========================================
        \node (head) [caja, below=2.5cm of backbone, fill=green!10] {YOLO Head \\ (Capas Densas)};
        
        \node (box) [caja, below left=1cm and 0.5cm of head] {Regresión Bounding Box \\ $(x, y, w, h)$};
        \node (obj) [caja, below=1cm of head] {Puntaje de Confianza \\ (\textit{Objectness})};
        \node (cls) [caja, below right=1cm and 0.5cm of head] {Clasificación de Clase \\ Probabilidad (A-Z)};

        \node (nms) [caja, below=1.5cm of obj] {Filtro NMS \\ (Supresión de No Máximos)};
        \node (eval) [matriz, right=1.5cm of nms] {Evaluación mAP};
        \node (final) [caja, below=1cm of nms, fill=red!10, thick] {Predicción Final de la Seña};

        \draw [flecha] (tensor_out.south) |- (head.east);
        
        \draw [flecha] (head.south) -- ++(0,-0.5) -| (box.north);
        \draw [flecha] (head.south) -- (obj.north);
        \draw [flecha] (head.south) -- ++(0,-0.5) -| (cls.north);

        \draw [flecha] (box.south) -- ++(0,-0.5) -| (nms.north);
        \draw [flecha] (obj.south) -- (nms.north);
        \draw [flecha] (cls.south) -- ++(0,-0.5) -| (nms.north);
        
        \draw [flecha] (nms) -- (final);
        \draw [flecha] (nms) -- (eval);

        \begin{scope}[on background layer]
            \node[draw=black, thick, dotted, fill=gray!5, inner sep=15pt, fit=(head) (box) (obj) (cls) (nms) (eval) (final)] (task_layer) {};
            \node[anchor=south east, text=black, font=\bfseries\footnotesize] at (task_layer.south east) {Task Layers (YOLO Heads)};
        \end{scope}

    \end{tikzpicture}
    }
    \caption{Diagrama de arquitectura profunda basado en YOLO para la detección y clasificación simultánea de la LSC, estructurado en capas de pre-procesamiento, extracción compartida y tareas de predicción.}
    \label{fig:arquitectura_yolo_compleja}
\end{figure}

\end{document}
